% Annexe B — Extraits de code et de configuration commentés
\label{chap:annexe-code}

Cette annexe rassemble les extraits de code et les éléments de configuration les plus représentatifs du projet. Chacun renvoie à la partie du dossier qui l'exploite.

\section*{Composition Docker de l'environnement local}
\phantomsection\label{ann:docker-local}
Schéma simplifié des services orchestrés par \texttt{docker-compose.yml} en développement, référencé à la \hyperref[subsec:env-local]{partie~\ref*{subsec:env-local}}.

\begin{figure}[htbp]
    \centering
    \includegraphics[width=0.66\linewidth]{dockercomposelocal-simplifié}
    \caption{Environnement Docker local simplifié de BrackX.}
\end{figure}

\noindent\hyperref[subsec:env-local]{$\hookleftarrow$ Retour à la partie~\ref*{subsec:env-local} (Environnement de développement local).}

\FloatBarrier
\section*{Pipeline d'intégration continue (\texttt{.gitlab-ci.yml})}
\phantomsection\label{ann:ci-extract}
Extrait du modèle commun \texttt{.ci} (règle de déclenchement, image, couverture) et du job \texttt{ci:backend}, suivi d'une exécution réussie du pipeline. Éléments référencés à la \hyperref[subsec:ci-pipeline]{partie~\ref*{subsec:ci-pipeline}}.

\begin{figure}[htbp]
    \centering
    \includegraphics[width=0.8\linewidth]{cigitlab-simplifié}
    \caption{Extrait commenté du pipeline d'intégration continue (\texttt{.gitlab-ci.yml}).}
\end{figure}

\begin{figure}[htbp]
    \centering
    \includegraphics[width=0.82\linewidth]{cireussi}
    \caption{Exécution réussie du pipeline~: couverture (82,81~\% de statements), 208 tests passés et jobs \texttt{ci:backend}, \texttt{ci:frontend} et \texttt{ci:docker-build} au vert.}
\end{figure}

\noindent\hyperref[subsec:ci-pipeline]{$\hookleftarrow$ Retour à la partie~\ref*{subsec:ci-pipeline} (Description du pipeline).}

\FloatBarrier
\section*{Extraits de code du cœur métier}
\phantomsection\label{ann:code-metier}
Briques les plus représentatives de la logique du module tournois, exploitées au chapitre~2 et testées comme le décrit la partie~\ref*{sec:tests-unitaires}. Le contrat commun des formats est défini par l'interface \texttt{GenerateurFormat}, que chaque format implémente, une fabrique associant le type de phase à la bonne classe. La logique de classement départage les égalités par confrontation directe puis par une cascade scalaire.

\begin{figure}[htbp]
    \centering
    \includegraphics[width=0.47\linewidth]{tournoi-interface}
    \caption{Interface \texttt{GenerateurFormat}, contrat commun à tous les formats.}
\end{figure}

\begin{figure}[htbp]
    \centering
    \includegraphics[width=0.47\linewidth]{tournoi-generateur-roundrobin}
    \caption{Générateur du format poules (round-robin) implémentant l'interface.}
\end{figure}

\begin{figure}[htbp]
    \centering
    \includegraphics[width=0.68\linewidth]{tournoi-generateur}
    \caption{Fabrique associant chaque type de phase à son générateur (\texttt{generateur-format.factory}).}
\end{figure}

\begin{figure}[htbp]
    \centering
    \includegraphics[width=0.47\linewidth]{classement-utils}
    \caption{Cascade de départage des classements (\texttt{classement.util}).}
\end{figure}

\FloatBarrier
\section*{Schéma de données et migration Prisma}
\phantomsection\label{ann:prisma}
Extrait du \texttt{schema.prisma} (approche schéma-first) et d'une migration versionnée, en complément du chapitre~2.

\begin{figure}[htbp]
    \centering
    \includegraphics[width=0.44\linewidth]{prisma}
    \caption{Extrait du schéma Prisma décrivant une partie des entités du domaine.}
\end{figure}

\begin{figure}[htbp]
    \centering
    \includegraphics[width=0.38\linewidth]{prisma-migration}
    \caption{Exemple de migration Prisma versionnée dans le dépôt.}
\end{figure}
